% BANKS-SOURCE: pdf=46 print=36 kind=solution-section id=chapter03-numbered-solutions
\begin{bankssolutions}{Solutions to Problems for Chapter 3}

% BANKS-SOURCE: pdf=46 print=36 kind=solution id=3.1
\BanksSolution{3.1}
Choose the ordering $t_1>t_2>\cdots>t_n$.  Inserting the identity
$\sum_r\ket r\bra r=1$ between successive operators gives
\begin{align*}
 G_n(t_1,\ldots,t_n)
 &\equiv \langle0|x(t_1)\cdots x(t_n)|0\rangle \\
 &=\sum_{r_1,\ldots,r_{n-1}}
 \langle0|x|r_1\rangle\langle r_1|x|r_2\rangle\cdots
 \langle r_{n-1}|x|0\rangle\\[-2pt]
 &\qquad\times
 \exp\!\left[\begin{aligned}
 &-\ii E_{r_1}(t_1-t_2)-\ii E_{r_2}(t_2-t_3)-\cdots\\[-2pt]
 &{}-\ii E_{r_{n-1}}(t_{n-1}-t_n)
 \end{aligned}\right],
\end{align*}
where $H\ket r=E_r\ket r$ and $E_0=0$.  In particular,
\begin{equation*}
 G_2(\tau)=\sum_r |\langle0|x|r\rangle|^2\ee^{-\ii E_r\tau},
 \qquad \tau>0.
\end{equation*}
The Fourier transform, or equivalently the Laplace transform after a small
positive damping of $\tau$, has delta-function support at every $E_r$ for
which $\langle0|x|r\rangle\ne0$.  The higher functions contain the transition
matrix elements and hence reveal states that first occur after repeated
applications of $x$.  Thus the full set of Green functions determines the
spectrum in the cyclic subspace generated by $x\ket0$.  Under the usual
assumption in this reconstruction problem that this subspace is the whole
Hilbert space, it determines every eigenvalue of $H$.  Degeneracies and matrix
elements are read from the coefficients of the corresponding spectral terms.

Time derivatives of a time-ordered function insert $\dot x$.  Taking all time
arguments to a common value through a fixed ordered side therefore gives
expectation values of ordered products of $x$ and $\dot x$.  Contact terms at
coincident times are known canonical commutators, so they can be subtracted.
For a Hamiltonian with kinetic term $p^2/(2M)$, $p=M\dot x$; consequently all
moments $\langle0|x^m p^n|0\rangle$ are determined.  A convenient equivalent
package is the Weyl characteristic function
\begin{equation*}
 \chi(a,b)=\left\langle0\left|\exp\!\bigl(\ii(ax+bp)\bigr)\right|0\right\rangle.
\end{equation*}
Its Taylor coefficients are the completely symmetrized moments, which are
obtained from the ordered moments by repeated use of $[x,p]=\ii$.  Fourier
inversion of $\chi$ gives the Wigner distribution and hence the pure-state
density kernel
\begin{equation*}
 \rho_0(x,y)=\psi_0(x)\psi_0(y)^*.
\end{equation*}
Factoring this rank-one kernel determines $\psi_0$ up to one constant phase.
This step assumes the standard determinate moment or analytic-characteristic
function condition; the Green functions themselves provide the characteristic
function when moments are not used as a formal series.

With the ground-state energy set to zero, the Schrödinger equation is
\begin{equation*}
 \left[-\frac{1}{2M}\frac{\dd^2}{\dd x^2}+V(x)\right]\psi_0(x)=0.
\end{equation*}
At every point where $\psi_0(x)\ne0$ this gives
\begin{equation*}
 V(x)=\frac{1}{2M}\frac{\psi_0''(x)}{\psi_0(x)}.
\end{equation*}
For a regular one-dimensional confining potential the ground state can be
chosen nodeless, and the expression extends by continuity.  An arbitrary
ground-state energy would add that constant to the right-hand side.  Therefore
the Green functions determine both the Hamiltonian (up to the stated energy and
phase conventions) and its spectrum.

% BANKS-SOURCE: pdf=46 print=36 kind=solution id=3.2
\BanksSolution{3.2}
Write the spatial source as
\begin{equation*}
 \rho(\mathbf x)=q_1\delta^{D-1}(\mathbf x)+q_2\delta^{D-1}(\mathbf x-\mathbf R),
 \qquad J(t,\mathbf x)=\rho(\mathbf x)\,\mathbf 1_{[-T,T]}(t).
\end{equation*}
It is simplest to evaluate the Gaussian in Euclidean signature and then
continue back to the Lorentzian vacuum amplitude.  For the free scalar,
\begin{equation*}
 Z_E[J]=\exp\!\left[-\frac12\int\frac{\dd^D p_E}{(2\pi)^D}
 \frac{|\widetilde J(p_E)|^2}{p_E^2+m^2}\right],
 \qquad
 \widetilde J(p_E^0,\mathbf p)=\frac{2\sin(p_E^0T)}{p_E^0}\widetilde\rho(\mathbf p).
\end{equation*}
As $T\to\infty$,
\begin{equation*}
 \left(\frac{2\sin(p_E^0T)}{p_E^0}\right)^2
 \longrightarrow 4\pi T\,\delta(p_E^0)
\end{equation*}
in the sense of distributions.  Hence
\begin{align*}
 \log Z_E[J]
 &=-T\int\frac{\dd^{D-1}\mathbf p}{(2\pi)^{D-1}}
 \frac{|q_1+q_2\ee^{-\ii\mathbf p\cdot\mathbf R}|^2}{\mathbf p^2+m^2}+o(T)\\
 &=-T\left[(q_1^2+q_2^2)I_D(0)+2q_1q_2I_D(R)\right]+o(T),
\end{align*}
where
\begin{equation*}
\begin{aligned}
 I_D(R)&=\int\frac{\dd^{D-1}\mathbf p}{(2\pi)^{D-1}}
 \frac{\ee^{+\ii\mathbf p\mathbin{\cdot}\mathbf R}}{\mathbf p^2+m^2}\\
 &=\frac{m^{(D-3)/2}}{(2\pi)^{(D-1)/2}R^{(D-3)/2}}
 K_{(D-3)/2}(mR),\\
 I_4(R)&=\frac{\ee^{-mR}}{4\pi R}.
\end{aligned}
\end{equation*}
The value $I_D(0)$ depends on the ultraviolet regulator and is independent of
$R$.  The same large-$T$ limit has a spectral interpretation.  The Euclidean
evolution with the source turned on is generated by the Hamiltonian with the
static coupling $-\int\rho\phi$.  Its vacuum matrix element has the form
\begin{equation*}
 \langle0|\exp(-2T H_\rho)|0\rangle
 =C\ee^{-2V(R)T}\bigl(1+o(1)\bigr),
\end{equation*}
because the lowest-energy state with the two prescribed localized
disturbances dominates.  Comparing exponents gives
\begin{equation*}
 V(R)=\frac12(q_1^2+q_2^2)I_D(0)+q_1q_2 I_D(R).
\end{equation*}
The first term is a sum of source self-energies and can be absorbed into the
additive constant requested in the problem.  Thus, with the sign convention
for the source used here,
\begin{equation*}
 V_{\mathrm{int}}^{(D)}(R)=q_1q_2 I_D(R),
 \qquad
 V_{\mathrm{int}}^{(4)}(R)=q_1q_2\frac{\ee^{-mR}}{4\pi R}.
\end{equation*}
The Lorentzian result follows from the contour rotation $p^0\to\ii p_E^0$
with $t=-\ii\tau$ and retained spatial coordinates;
the Yukawa factor is the zero-frequency massive propagator.  A single virtual
scalar exchange therefore gives the force between the static disturbances.

% BANKS-SOURCE: pdf=46 print=36 kind=solution id=3.3
\BanksSolution{3.3}
Let $t=x^0-y^0$ and abbreviate
\begin{equation*}
 A(t)=\langle0|\phi(x)\phi(y)|0\rangle,
 \qquad B(t)=\langle0|\phi(y)\phi(x)|0\rangle.
\end{equation*}
For two bosonic fields,
\begin{equation*}
 G(x,y)=\langle0|T\phi(x)\phi(y)|0\rangle
 =\theta(t)A(t)+\theta(-t)B(t).
\end{equation*}
Differentiating once gives
\begin{equation*}
 \partial_0G=\langle0|T\dot\phi(x)\phi(y)|0\rangle
 +\delta(t)\langle0|[\phi(x),\phi(y)]|0\rangle.
\end{equation*}
The equal-time canonical relation makes the last commutator zero, but it must
be retained while taking the second derivative.  Direct differentiation before
using the equal-time relation gives
\begin{align*}
 \partial_0^2G={}&\langle0|T\ddot\phi(x)\phi(y)|0\rangle
 +\delta'(t)\langle0|[\phi(x),\phi(y)]|0\rangle\\
 &+2\delta(t)\langle0|[\dot\phi(x),\phi(y)]|0\rangle.
\end{align*}
If $C(t)=\langle0|[\phi(x),\phi(y)]|0\rangle$, then
$C(0)=0$ and $C'(0)=\langle0|[\dot\phi(x),\phi(y)]|0\rangle$.
The distribution identity $\delta'(t)C(t)=-\delta(t)C'(0)$ therefore leaves
one contact term:
\begin{equation*}
 \partial_0^2G=\langle0|T\ddot\phi(x)\phi(y)|0\rangle
 +\delta(t)\langle0|[\dot\phi(x),\phi(y)]|0\rangle.
\end{equation*}
With $\pi=\dot\phi$ and $[\phi(t,\mathbf x),\pi(t,\mathbf y)]
=\ii\delta^{D-1}(\mathbf x-\mathbf y)$,
\begin{equation*}
 \langle0|[\dot\phi(x),\phi(y)]|0\rangle
 =-\ii\delta^{D-1}(\mathbf x-\mathbf y),
\end{equation*}
so
\begin{equation*}
 \partial_0^2G=\langle0|T\ddot\phi(x)\phi(y)|0\rangle
 -\ii\delta^D(x-y).
\end{equation*}

For $n$ further fields, differentiating every Heaviside factor produces the
same contact term for each possible $y_j$.  The result is
\begin{align*}
 \partial_0^2\langle0|T\phi(x)\phi(y_1)\cdots\phi(y_n)|0\rangle
={}&\langle0|T\ddot\phi(x)\phi(y_1)\cdots\phi(y_n)|0\rangle\\
 &-\ii\sum_{j=1}^n\delta^D(x-y_j)
 \\[-2pt]
 &\quad{}\times\langle0|T\phi(y_1)\cdots\widehat{\phi(y_j)}\cdots\phi(y_n)|0\rangle.
\end{align*}
The hat means that the $j$th field is omitted.  The spatial part of $\partial^2$
acts without differentiating time-ordering functions.  Since
$\Box=-\partial_0^2+\nabla^2$, the contact term in the full d'Alembertian is
$+\ii\delta^D$.  Using the Heisenberg equation $\Box\phi(x)=V'(\phi(x))$
gives the operator insertion $V'(\delta/(\ii\delta J))$.  We adopt throughout this solution the source
convention printed in (3.3),
$Z[J]=\langle T\exp(+\ii\int\dd^D x\,J(x)\phi(x))\rangle$.  In the $n$th term of
$\delta Z/(\ii\delta J(x))$, the $n$ contact terms in the preceding
equation carry the factor
$(-\ii)\ii^n n/n!=\ii^{n-1}/(n-1)!$.  The minus sign in
$\Box=-\partial_0^2+\nabla^2$ reverses this contact contribution.  After the remaining
source integrations, the full d'Alembertian contact term sums to $-J(x)Z[J]$.
Thus the contact contribution appears as the $-J(x)$ in the
Schwinger--Dyson bracket:
\begin{equation*}
 \Box\frac{\delta Z}{\ii\delta J(x)}
 =\left[V'\!\left(\frac{\delta}{\ii\delta J(x)}\right)-J(x)\right]Z[J].
\end{equation*}
The contact term is the origin of the source term; omitting the derivative of
the Heaviside functions would incorrectly give a homogeneous equation.  The
The historical source's printed $+J$ in this bracket is consequently a sign error relative to
its own $+\ii\int J\phi$ definition.

% BANKS-SOURCE: pdf=46 print=36 kind=solution id=3.4
\BanksSolution{3.4}
After Wick rotation, the free and interaction parts are
\begin{equation*}
 S_E=\int\dd^D x_E\left[\frac12(\partial_E\phi)^2
 +\frac12m_0^2\phi^2+\frac{\lambda_3}{3!}\phi^3
 +\frac{\lambda_4}{4!}\phi^4\right].
\end{equation*}
The Euclidean propagator and the two scalar vertices are
\begin{equation*}
 \widetilde D_E(k_E)=\frac{1}{k_E^2+m_0^2},
 \qquad V_3=-\lambda_3,
 \qquad V_4=-\lambda_4.
\end{equation*}
The connected tree-level four-point function has one quartic contact graph and
three graphs made from two cubic vertices.  With all four external propagators
shown explicitly,
\begin{align*}
 \widetilde G_{4,E}^{\rm tree}(k_{E,1},\ldots,k_{E,4})
 &=(2\pi)^D\delta^D\!\left(\sum_r k_{E,r}\right)
 \prod_{r=1}^4\frac{1}{k_{E,r}^2+m_0^2}\,\mathcal A_E,\\
 \mathcal A_E
 &=-\lambda_4+\frac{\lambda_3^2}{(k_{E,1}+k_{E,2})^2+m_0^2}\\[-2pt]
 &\quad+\frac{\lambda_3^2}{(k_{E,1}-k_{E,3})^2+m_0^2}\\[-2pt]
 &\quad+\frac{\lambda_3^2}{(k_{E,1}-k_{E,4})^2+m_0^2}.
\end{align*}
The three denominators are the $s$, $t$, and $u$ channels after continuation.
There are no other connected four-point trees: a tree built only from one
quartic vertex is the contact graph, and a tree with two cubic vertices has one
internal line and gives one of the three partitions of four labeled external
legs into two pairs.
The full, unconnected four-point function also contains the three free Wick
pairings.  Those pairings describe independent propagation and are removed
when the connected function is used in the LSZ formula.

The continuation $p^0\to\ii p_E^0$ gives $k_E^2=p^2$ and
\begin{equation*}
 \widetilde D_F(p)=\frac{-\ii}{p^2+m_0^2-\ii\varepsilon},
 \qquad V_3=-\ii\lambda_3,
 \qquad V_4=-\ii\lambda_4.
\end{equation*}
Thus the amputated Lorentzian amplitude is
\begin{equation*}
 \ii\mathcal M(p_1,p_2\to p_3,p_4)
 =-\ii\lambda_4
 +\frac{(-\ii\lambda_3)^2\ii}{s-m_0^2+\ii\varepsilon}
 +\frac{(-\ii\lambda_3)^2\ii}{t-m_0^2+\ii\varepsilon}
 +\frac{(-\ii\lambda_3)^2\ii}{u-m_0^2+\ii\varepsilon},
\end{equation*}
or, away from the poles,
\begin{equation*}
 \mathcal M=-\lambda_4-\frac{\lambda_3^2}{s-m_0^2}
 -\frac{\lambda_3^2}{t-m_0^2}-\frac{\lambda_3^2}{u-m_0^2}.
\end{equation*}

In the semiclassical normalization of the text, a $k$-point interaction is
$O(g^{k-2})$.  Therefore $\lambda_3=O(g)$ and $\lambda_4=O(g^2)$; each term
in the four-point tree amplitude is $O(g^2)$.  At tree level the two-point
function is the free propagator, so its pole is at $p^2=-m_0^2$ and its residue
is one.  Hence the tree values are $m=m_0$ and $Z=1$.  A self-energy or a
momentum-dependent residue first requires a loop.

The LSZ formula removes the four external propagators and multiplies by
$Z^{-1/2}$ for each leg.  Since $Z=1$, it gives exactly the amputated
expression above.  On shell, $p_r^2=-m_0^2$ and
$p_1+p_2=p_3+p_4$.  Expanding the definitions gives
\begin{align*}
 s+t+u
 &=-\left[3p_1^2+p_2^2+p_3^2+p_4^2
 +2p_1\!\cdot(p_2-p_3-p_4)\right]\\
 &=-\left[-6m_0^2+2m_0^2\right]=4m_0^2.
\end{align*}
In the center-of-mass frame set
\begin{equation*}
 \begin{aligned}
 p_1&=(E,\mathbf p),\quad p_2=(E,-\mathbf p),\\[-2pt]
 p_3&=(E,\mathbf p'),\quad p_4=(E,-\mathbf p'),\\[-2pt]
 &\qquad |\mathbf p|=|\mathbf p'|=p,
 \end{aligned}
\end{equation*}
and let $\theta$ be the angle between $\mathbf p$ and $\mathbf p'$.  Then
\begin{equation*}
 \begin{aligned}
 s&=4E^2,\\[-2pt]
 t&=-2p^2(1-\cos\theta)=-4p^2\sin^2\frac\theta2,\\[-2pt]
 u&=-2p^2(1+\cos\theta)=-4p^2\cos^2\frac\theta2,
 \end{aligned}
\end{equation*}
where $E^2=\mathbf p^2+m_0^2$.  Equivalently,
\begin{equation*}
 t=-\frac{s-4m_0^2}{2}(1-\cos\theta),
 \qquad u=-\frac{s-4m_0^2}{2}(1+\cos\theta).
\end{equation*}
The real scalar has identical external particles, so relabeling the legs
permutes $s,t,u$ and leaves the amplitude unchanged.  The same statement
relates the physical channel regions by analytic continuation of external
momenta.  This is the tree-level crossing relation.

% BANKS-SOURCE: pdf=47 print=37 kind=solution id=3.5
\BanksSolution{3.5}
Let $a$ be the number of cubic vertices, $b$ the number of quartic vertices,
$I$ the number of internal lines, and $E$ the number of labeled external legs.
Every connected graph obeys
\begin{equation*}
 3a+4b=2I+E,
 \qquad L=I-a-b+1,
 \qquad E=a+2b+2-2L.
\end{equation*}
The complete vertex-count list for $E\le8$ is
\begin{center}
\small
\begin{tabular}{c|c|l}
 $L$ & $E$ & (cubic, quartic) $(a,b)$ \\ \hline
 $0$ & $2$ & $(0,0)$ \\
 $0$ & $3$ & $(1,0)$ \\
 $0$ & $4$ & $(0,1),(2,0)$ \\
 $0$ & $5$ & $(1,1),(3,0)$ \\
 $0$ & $6$ & $(0,2),(2,1),(4,0)$ \\
 $0$ & $7$ & $(1,2),(3,1),(5,0)$ \\
 $0$ & $8$ & $(0,3),(2,2),(4,1),(6,0)$ \\ \hline
 $1$ & $1$ & $(1,0)$ \\
 $1$ & $2$ & $(0,1),(2,0)$ \\
 $1$ & $3$ & $(1,1),(3,0)$ \\
 $1$ & $4$ & $(0,2),(2,1),(4,0)$ \\
 $1$ & $5$ & $(1,2),(3,1),(5,0)$ \\
 $1$ & $6$ & $(0,3),(2,2),(4,1),(6,0)$ \\
 $1$ & $7$ & $(1,3),(3,2),(5,1),(7,0)$ \\
 $1$ & $8$ & $(0,4),(2,3),(4,2),(6,1),(8,0)$ \\ \hline
 $2$ & $0$ & $(0,1),(2,0)$ \\
 $2$ & $1$ & $(1,1),(3,0)$ \\
 $2$ & $2$ & $(0,2),(2,1),(4,0)$ \\
 $2$ & $3$ & $(1,2),(3,1),(5,0)$ \\
 $2$ & $4$ & $(0,3),(2,2),(4,1),(6,0)$ \\
 $2$ & $5$ & $(1,3),(3,2),(5,1),(7,0)$ \\
 $2$ & $6$ & $(0,4),(2,3),(4,2),(6,1),(8,0)$ \\
 $2$ & $7$ & $(1,4),(3,3),(5,2),(7,1),(9,0)$ \\
 $2$ & $8$ & $(0,5),(2,4),(4,3),(6,2),(8,1),(10,0)$
\end{tabular}
\end{center}
The first tree row contains the free propagator; all other entries are
interaction graphs.  The empty normalization graph is not a connected graph;
self-contractions are retained in the enumeration below.

Represent a topology by vertices $v_i$ of valence $k_i=3$ or $4$, external-leg
counts $e_i$, multiplicities $n_{ij}$ of lines between distinct vertices, and
$\ell_i$ self-contracted lines at $v_i$.  They satisfy
\begin{equation*}
 e_i+\sum_{j\ne i}n_{ij}+2\ell_i=k_i,
 \qquad \sum_i e_i=E,
 \qquad I=\sum_{i<j}n_{ij}+\sum_i\ell_i.
\end{equation*}
The symmetry factor of the graph with its external labels fixed is
\begin{equation*}
 S_G=\frac{1}{|\operatorname{Aut}G|},
\end{equation*}
where automorphisms fix every external label.  More explicitly, if a
permutation of internal vertices preserves their valences, external incidence,
and all $n_{ij}$ and $\ell_i$, then it contributes; in addition every bundle of
$n_{ij}$ parallel lines contributes a factor $n_{ij}!$, and every self-loop
contributes its reversal and its permutation factors $2^{\ell_i}\ell_i!$.
Thus one may compute
\begin{equation*}
 |\operatorname{Aut}G|
 =|\operatorname{Aut}_{\rm vert}G|
 \prod_{i<j}n_{ij}!\prod_i2^{\ell_i}\ell_i!.
\end{equation*}
This is the geometric count, and it is also the Wick count because
\begin{equation*}
 S_G=\frac{N_G}{a!\,b!\,(3!)^a(4!)^b},
\end{equation*}
where $N_G$ is the number of contractions producing the fixed labeled graph.

To make the answer checkable graph by graph, use the following canonical
incidence record.  The first $a$ vertices have $k_i=3$ and the next $b$ have
$k_i=4$.  The vector $e=(e_1,\ldots,e_V)$ gives the external-leg blocks in
that order.  The vector $n$ lists
$(n_{12},n_{13},\ldots,n_{1V},n_{23},\ldots,n_{V-1,V})$, and $\ell$ lists
the self-contractions $(\ell_1,\ldots,\ell_V)$.  For the compact
representative, external-block counts are listed in nondecreasing order within
each valence class.  After this choice the external blocks are marked in their
displayed order.  Thus a topology is canonical up to permutations of
unmarked vertices of the same valence; permutations of the external labels
give channels of the same record and the same factor.

The complete per-topology enumeration is generated by the self-contained code
below.  It uses only the Python standard library, fixes the traversal order,
prints one record containing $(e,n,\ell,|\operatorname{Aut}G|,S_G,N_G)$ for
each topology, and finishes with a row count and SHA-256 digest of the complete
output.  The default bounds are exactly $E\le8$ and $L\le2$; the environment
variables shown in the first two lines are available only for a smaller smoke
run.  The code is embedded here so the exhaustive output is reproducible from
this solution file without an external topology list.

\begingroup
\small
\begin{verbatim}
from collections import defaultdict
from fractions import Fraction
from hashlib import sha256
from itertools import combinations_with_replacement, permutations
from itertools import product
from math import factorial
import os

MAX_E = int(os.environ.get("BANKS_MAX_E", "8"))
MAX_L = int(os.environ.get("BANKS_MAX_L", "2"))

def bounded_compositions(total, caps):
    out = []
    suffix = [0] * (len(caps) + 1)
    for i in range(len(caps) - 1, -1, -1):
        suffix[i] = suffix[i + 1] + caps[i]
    def rec(i, left, values):
        if i == len(caps):
            if left == 0:
                out.append(tuple(values))
            return
        lo = max(0, left - suffix[i + 1])
        hi = min(caps[i], left)
        for value in range(lo, hi + 1):
            values.append(value)
            rec(i + 1, left - value, values)
            values.pop()
    rec(0, total, [])
    return out

def multigraphs(degrees):
    v = len(degrees)
    remaining = list(degrees)
    matrix = [[0] * v for _ in range(v)]
    def rec(i):
        if i == v:
            if not any(remaining):
                yield tuple(tuple(row) for row in matrix)
            return
        for loops in range(remaining[i] // 2 + 1):
            remaining[i] -= 2 * loops
            matrix[i][i] = loops
            need = remaining[i]
            if i == v - 1:
                if need == 0:
                    yield from rec(i + 1)
            else:
                for values in bounded_compositions(
                        need, remaining[i + 1:]):
                    for j, value in enumerate(values, i + 1):
                        matrix[i][j] = matrix[j][i] = value
                        remaining[j] -= value
                    remaining[i] = 0
                    yield from rec(i + 1)
                    remaining[i] = need
                    for j, value in enumerate(values, i + 1):
                        matrix[i][j] = matrix[j][i] = 0
                        remaining[j] += value
            matrix[i][i] = 0
            remaining[i] += 2 * loops
    yield from rec(0)

def connected(matrix):
    if len(matrix) == 0:
        return False
    if len(matrix) == 1:
        return True
    seen = {0}
    stack = [0]
    while stack:
        i = stack.pop()
        for j, multiplicity in enumerate(matrix[i]):
            if i != j and multiplicity and j not in seen:
                seen.add(j)
                stack.append(j)
    return len(seen) == len(matrix)

def incidence_key(matrix, permutation):
    v = len(matrix)
    return tuple(matrix[permutation[i]][permutation[j]]
                 for i in range(v) for j in range(i, v))

def group_permutations(groups, v):
    if not groups:
        return [tuple(range(v))]
    choices = [list(permutations(group)) for group in groups]
    out = []
    for selected in product(*choices):
        permutation = list(range(v))
        for group, image in zip(groups, selected):
            for old, new in zip(group, image):
                permutation[old] = new
        out.append(tuple(permutation))
    return out

def canonical(matrix, allowed_permutations):
    return min(incidence_key(matrix, p) for p in allowed_permutations)

def equal_under(key, permutation):
    v = len(permutation)
    def index(i, j):
        if i > j:
            i, j = j, i
        return i * v - i * (i - 1) // 2 + (j - i)
    return all(key[index(i, j)] ==
               key[index(permutation[i], permutation[j])]
               for i in range(v) for j in range(i, v))

def incidence(key, v):
    def index(i, j):
        if i > j:
            i, j = j, i
        return i * v - i * (i - 1) // 2 + (j - i)
    n = tuple(key[index(i, j)] for i in range(v)
              for j in range(i + 1, v))
    ell = tuple(key[index(i, i)] for i in range(v))
    return n, ell

digest = sha256()
rows = 0
def emit(l, e, a, b, external, n, ell, aut, symmetry, contractions):
    global rows
    line = (f"L={l} E={e} a={a} b={b} e={external} n={n} "
            f"ell={ell} Aut={aut} S={symmetry} N={contractions}")
    print(line)
    digest.update((line + "\\n").encode())
    rows += 1

emit(0, 2, 0, 0, (), (), (), 1, Fraction(1), 1)
for l in range(MAX_L + 1):
    for e in range(MAX_E + 1):
        if l == 0 and e == 2:
            continue
        target = e - 2 + 2 * l
        if target < 0:
            continue
        for b in range(target // 2 + 1):
            a = target - 2 * b
            v = a + b
            kinds = (3,) * a + (4,) * b
            for cubic_e in combinations_with_replacement(range(4), a):
                for quartic_e in combinations_with_replacement(
                        range(5), b):
                    external = cubic_e + quartic_e
                    if sum(external) != e:
                        continue
                    degrees = tuple(k - x
                                    for k, x in zip(kinds, external))
                    if sum(degrees) % 2:
                        continue
                    zero_groups = []
                    for kind in (3, 4):
                        group = [i for i, k in enumerate(kinds)
                                 if k == kind and external[i] == 0]
                        if group:
                            zero_groups.append(group)
                    allowed_permutations = group_permutations(
                        zero_groups, v)
                    seen = set()
                    for matrix in multigraphs(degrees):
                        if not connected(matrix):
                            continue
                        key = canonical(matrix, allowed_permutations)
                        if key in seen:
                            continue
                        seen.add(key)
                        aut_vertices = sum(
                            equal_under(key, p)
                            for p in allowed_permutations
                        )
                        n, ell = incidence(key, v)
                        aut = aut_vertices
                        for value in n:
                            aut *= factorial(value)
                        for value in ell:
                            aut *= (2 ** value) * factorial(value)
                        denominator = (factorial(a) * factorial(b) *
                                      factorial(3) ** a *
                                      factorial(4) ** b)
                        assert denominator % aut == 0
                        emit(l, e, a, b, external, n, ell, aut,
                             Fraction(1, aut), denominator // aut)
print(f"ROWS={rows}")
print(f"SHA256={digest.hexdigest()}")
\end{verbatim}
\endgroup

The first records include the free propagator
$(L,E,a,b)=(0,2,0,0)$ with $e=n=\ell=()$, $|\operatorname{Aut}G|=1$,
$S_G=1$, and $N_G=1$.  A cubic three-point vertex has
$(e,n,\ell)=( (3),(),(0) )$, $|\operatorname{Aut}G|=1$, $S_G=1$, and
$N_G=3!=6$.  The quartic two-point tadpole has
$(e,n,\ell)=((2),(),(1))$, $|\operatorname{Aut}G|=2$, $S_G=1/2$, and
$N_G=12$.  The one-loop cubic two-point bubble has
$(e,n,\ell)=((1,1),(2),(0,0))$, $|\operatorname{Aut}G|=2$, $S_G=1/2$,
and $N_G=36$.  The two-loop quartic sunset has
$(e,n,\ell)=((1,1),(3),(0,0))$, $|\operatorname{Aut}G|=6$, $S_G=1/6$,
and $N_G=192$.

The three-parallel-line cubic--cubic vacuum theta graph is the valid
unmarked example: $a=2$, $b=0$, $E=0$, $e=(0,0)$, $n=(3)$, and
$\ell=(0,0)$.  The line permutations and the interchange of the two cubic
vertices give $|\operatorname{Aut}G|=3!\,2=12$, hence $S_G=1/12$ and
$N_G=2!(3!)^2/12=6$.  For a marked vertex, use instead the valid
$L=2$, $E=1$, $(a,b)=(1,1)$ graph with
$e=(0,1)$, $n=(3)$, and $\ell=(0,0)$: the external leg marks the quartic
vertex, so $|\operatorname{Aut}G|=3!=6$, $S_G=1/6$, and
$N_G=3!\,4!/6=24$.  These records illustrate the same two calculations used
for every emitted row: geometrical automorphisms and Wick contractions.

% BANKS-SOURCE: pdf=47 print=37 kind=solution id=3.6
\BanksSolution{3.6}
Expand about a solution of the source-dependent classical equation by writing
\begin{equation*}
 \phi=\phi_{\rm cl}+g\chi.
\end{equation*}
The linear term vanishes because $\phi_{\rm cl}$ is stationary.  Taylor
expansion of the action gives
\begin{equation*}
 \frac1{g^2}S[\phi_{\rm cl}+g\chi]
 =\frac1{g^2}S[\phi_{\rm cl}]
 +\frac12\chi K\chi
 -\sum_{k\ge3}\frac{g^{k-2}}{k!}
 \int\dd^D x\,V^{(k)}(\phi_{\rm cl}(x))\chi(x)^k,
\end{equation*}
where $K$ is the quadratic fluctuation operator.  Signs from the potential
are included in the displayed coefficient and do not affect the power of $g$.
The propagator obtained from $K$ has power $g^0$, every $k$-point vertex has
power $g^{k-2}$, and an external line carries no additional power of $g$ in
the Green functions of $\chi$.

Let $V_k$ denote the number of $k$-vertices and put
\begin{equation*}
 V=\sum_{k\ge3}V_k,
 \qquad I=\text{number of internal lines}.
\end{equation*}
Counting half-edges gives
\begin{equation*}
 \sum_{k\ge3}kV_k=2I+E.
\end{equation*}
The power of $g$ in the diagram is therefore
\begin{align*}
 \sum_{k\ge3}(k-2)V_k
 &=\sum_k kV_k-2V\\
 &=2I+E-2V.
\end{align*}
For a connected graph, Euler's relation is $L=I-V+1$, so
\begin{equation*}
 2I+E-2V=2(I-V+1)-2+E=2L-2+E.
\end{equation*}
Thus every connected $E$-point graph has the uniform order
\begin{equation*}
 \boxed{g^{\,2L-2+E}}.
\end{equation*}
The result uses no factors on external lines.  If one instead asks for Green
functions of the original fluctuation $\delta\phi=g\chi$, each external
operator contributes an additional explicit $g$, as follows directly from the
decomposition; the convention in the problem is the $\chi$ convention.

\end{bankssolutions}
